\documentclass[12pt,a4paper]{article}

% ── Packages ──────────────────────────────────────────────────────────────────
\usepackage[utf8]{inputenc}
\usepackage[T1]{fontenc}
\usepackage[margin=2.5cm]{geometry}
\usepackage{amsmath}
\usepackage{amssymb}
\usepackage{graphicx}
\usepackage{booktabs}
\usepackage{longtable}
\usepackage{array}
\usepackage{hyperref}
\usepackage{xcolor}
\usepackage{listings}
\usepackage{caption}
\usepackage{float}
\usepackage{setspace}
\usepackage{titlesec}
\usepackage{enumitem}
\usepackage{tabularx}
\usepackage{multirow}

% ── Hyperref colours ──────────────────────────────────────────────────────────
\hypersetup{
  colorlinks=true,
  linkcolor=blue!60!black,
  citecolor=green!50!black,
  urlcolor=blue!70!black,
  pdftitle={DL4AI Final Project Report},
  pdfauthor={Student 230051}
}

% ── Code listings style ───────────────────────────────────────────────────────
\definecolor{codebg}{HTML}{F5F5F5}
\definecolor{codeframe}{HTML}{CCCCCC}
\definecolor{keyword}{HTML}{0000FF}
\definecolor{comment}{HTML}{008000}
\definecolor{string}{HTML}{BA2121}

\lstdefinestyle{python}{
  language=Python,
  backgroundcolor=\color{codebg},
  frame=single,
  rulecolor=\color{codeframe},
  basicstyle=\ttfamily\small,
  keywordstyle=\color{keyword}\bfseries,
  commentstyle=\color{comment}\itshape,
  stringstyle=\color{string},
  breaklines=true,
  breakatwhitespace=true,
  showstringspaces=false,
  tabsize=4,
  captionpos=b
}

\lstdefinestyle{bash}{
  language=bash,
  backgroundcolor=\color{codebg},
  frame=single,
  rulecolor=\color{codeframe},
  basicstyle=\ttfamily\small,
  keywordstyle=\color{keyword},
  breaklines=true,
  showstringspaces=false,
  captionpos=b
}

\lstdefinestyle{plain}{
  backgroundcolor=\color{codebg},
  frame=single,
  rulecolor=\color{codeframe},
  basicstyle=\ttfamily\small,
  breaklines=true,
  showstringspaces=false,
  captionpos=b
}

% ── Section title formatting ──────────────────────────────────────────────────
\titleformat{\section}{\large\bfseries}{\thesection.}{0.5em}{}[\titlerule]
\titleformat{\subsection}{\normalsize\bfseries}{\thesubsection.}{0.5em}{}
\titleformat{\subsubsection}{\normalsize\itshape}{\thesubsubsection.}{0.5em}{}

% ── Spacing ───────────────────────────────────────────────────────────────────
\setlength{\parskip}{6pt}
\setlength{\parindent}{0pt}
\onehalfspacing

% ── Document ──────────────────────────────────────────────────────────────────
\begin{document}

% ── Title Page ────────────────────────────────────────────────────────────────
\begin{titlepage}
  \centering
  \vspace*{2cm}
  {\LARGE\bfseries Deep Learning for Artificial Intelligence Final Project Report\par}
  \vspace{0.5cm}
  {\large Time-Series Data and Application to Stock Markets\par}
  \vspace{2cm}
  \begin{tabular}{ll}
    \textbf{Student ID:} & 230051 \\[4pt]
    \textbf{Course:}     & Deep Learning for Artificial Intelligence (DL4AI) \\[4pt]
    \textbf{Date:}       & May 2026 \\
  \end{tabular}
  \vfill
  \hrule
  \vspace{0.3cm}
\end{titlepage}

\tableofcontents
\newpage

% ══════════════════════════════════════════════════════════════════════════════
\section*{Abstract}
\addcontentsline{toc}{section}{Abstract}
% ══════════════════════════════════════════════════════════════════════════════

This project applies deep learning to financial time-series data from two markets:
Nasdaq (USA) and the Vietnam stock exchange (VNINDEX). A unified
\textbf{CNN-LSTM} (Convolutional Neural Network --- Long Short-Term Memory)
architecture is designed and applied across five tasks:
(1) multi-feature next-day price prediction on Nasdaq,
(2) price prediction with technical indicators on Vietnam data,
(3) binary trading signal classification (BUY / SELL),
(4) quantitative portfolio construction and evaluation, and
(5) production deployment as a full-stack service --- including a Flask REST API,
a multi-page \textbf{Streamlit SaaS dashboard}, and an \textbf{Apache Airflow
automation pipeline} backed by SQLite.
Experimental results show that the model achieves low MAPE ($\approx$1--2\,\%)
for short-horizon price regression, moderate signal classification performance
(AUC 0.46--0.53), both constructed portfolios outperform the market benchmark by
+5.6 percentage points during the 2022 evaluation period, and the daily pipeline
reliably generates 19 predictions across 11 tickers in under 60 seconds.

% ══════════════════════════════════════════════════════════════════════════════
\section{Introduction}
% ══════════════════════════════════════════════════════════════════════════════

Stock market forecasting is a long-standing and difficult problem in quantitative
finance due to the inherently noisy, non-stationary, and high-dimensional nature
of financial time-series data. Classical approaches such as ARIMA, GARCH, and
hand-crafted technical analysis rules have been largely superseded by machine
learning methods that can learn non-linear patterns from raw data.

Among deep learning models, \textbf{Long Short-Term Memory (LSTM)} networks are
naturally suited to sequential data but can struggle to efficiently extract local
short-term patterns from raw OHLCV features. \textbf{Convolutional Neural
Networks (CNN)}, while primarily known for image tasks, are effective at
extracting local temporal features in one dimension. Combining the two --- a
\textbf{CNN-LSTM} architecture --- leverages the strengths of both: CNN layers
extract local candlestick-level patterns, and the LSTM layer models long-range
temporal dependencies.

\subsection{Objectives}

This project addresses five interconnected tasks:

\begin{enumerate}[leftmargin=2em]
  \item \textbf{Nasdaq price forecasting} --- next-day price, $k$-th day ahead,
        and $k$ consecutive days multi-output.
  \item \textbf{Vietnam price forecasting} --- incorporating 12 technical
        indicators alongside OHLCV features.
  \item \textbf{Trading signal classification} --- predicting BUY and SELL
        signals from historical price features.
  \item \textbf{Portfolio management} --- constructing and back-testing two
        investor-profile portfolios (risk-taking vs.\ prudent) from the VNINDEX
        universe.
  \item \textbf{Deployment} --- REST API (Flask), interactive web SaaS
        (Streamlit), automated daily pipeline (Airflow + SQLite), and Docker
        containerisation.
\end{enumerate}

\subsection{Research Questions}

\begin{itemize}[leftmargin=2em]
  \item Can a CNN-LSTM model accurately predict short-horizon stock prices using
        only OHLCV features?
  \item Do technical indicators improve forecasting on Vietnam stock data?
  \item How reliably can a classifier detect BUY and SELL trading signals from
        historical price data alone?
  \item Does a model-driven signal strategy outperform a buy-and-hold baseline?
  \item How can a trained deep learning model be served as a production-ready
        REST API?
  \item Can the inference workflow be automated and operationalised as a daily
        data pipeline?
\end{itemize}

% ══════════════════════════════════════════════════════════════════════════════
\section{Datasets}
% ══════════════════════════════════════════════════════════════════════════════

\subsection{Nasdaq Dataset}

\begin{table}[H]
\centering
\caption{Nasdaq dataset properties}
\begin{tabularx}{\linewidth}{lX}
\toprule
\textbf{Property} & \textbf{Detail} \\
\midrule
Source       & \texttt{data\_nasdaq\_csv/csv/} --- one CSV per ticker \\
Coverage     & $\sim$2,500+ US-listed Nasdaq stocks, daily data \\
Date range   & Varies per ticker (AAPL: 1984--2016+) \\
Features     & \texttt{Low, Open, Volume, High, Close, Adjusted Close} (6 columns) \\
Target       & \texttt{Close} price \\
Demo ticker  & AAPL (Apple Inc.) --- longest complete history \\
Preprocessing& Sort by date, drop NaN rows, build 30-day sliding windows \\
\bottomrule
\end{tabularx}
\end{table}

\subsection{Vietnam Stock Dataset (VNINDEX)}

\begin{table}[H]
\centering
\caption{Vietnam dataset properties}
\begin{tabularx}{\linewidth}{lX}
\toprule
\textbf{Property} & \textbf{Detail} \\
\midrule
Source          & \texttt{data-vn-20230228/stock-historical-data/} \\
Filename format & \texttt{\{TICKER\}-\{EXCHANGE\}-History.csv} \\
Coverage        & 416 VNINDEX-listed tickers, data through February 2023 \\
Features        & \texttt{Open, High, Low, Close, Volume} (5 base) + 12 technical indicators = \textbf{17 features} \\
Target          & \texttt{Close} price \\
Demo ticker     & REE (Refrigeration Electrical Engineering Corporation) \\
Supplementary   & Quarterly financial ratios from \texttt{financial-ratio/} (D/E, ROE, ROA, P/E, P/B) \\
\bottomrule
\end{tabularx}
\end{table}

\subsection{Technical Indicators (12 features added to Vietnam data)}

\begin{table}[H]
\centering
\caption{Technical indicators computed for Vietnam data}
\begin{tabular}{llp{7cm}}
\toprule
\textbf{Category} & \textbf{Indicator} & \textbf{Description} \\
\midrule
Trend      & SMA-10, SMA-20        & 10-day and 20-day Simple Moving Average \\
Momentum   & MACD                  & EMA(12) $-$ EMA(26) \\
Momentum   & MACD Signal           & 9-day EMA of MACD \\
Momentum   & MACD Histogram        & MACD $-$ MACD Signal \\
Oscillator & RSI-14                & Relative Strength Index (14-day) \\
Volatility & BB Upper / BB Lower   & Bollinger Band boundaries (20-day, $\pm2\sigma$) \\
Volatility & BB Width              & BB Upper $-$ BB Lower \\
Volatility & BB \%B                & Position of price within bands \\
Volatility & ATR-14                & Average True Range (14-day) \\
Volume     & OBV                   & On-Balance Volume (cumulative signed volume) \\
\bottomrule
\end{tabular}
\end{table}

\subsection{Data Split Strategy}

All experiments use a \textbf{strictly chronological split} to prevent data
leakage:

\begin{table}[H]
\centering
\caption{Train/validation/test split}
\begin{tabular}{lll}
\toprule
\textbf{Partition} & \textbf{Ratio} & \textbf{Purpose} \\
\midrule
Train      & 70\,\% & Model parameter fitting \\
Validation & 15\,\% & Early stopping \& LR scheduling \\
Test       & 15\,\% & Final held-out evaluation \\
\bottomrule
\end{tabular}
\end{table}

For Section~4 portfolio management, an additional \textbf{252-trading-day
($\approx$1 year) hold-out} period is reserved at the end of each ticker's series
for out-of-sample performance simulation.

% ══════════════════════════════════════════════════════════════════════════════
\section{Methodology}
% ══════════════════════════════════════════════════════════════════════════════

\subsection{Sliding Window Construction}

All input sequences are constructed as fixed-length sliding windows:

\begin{itemize}[leftmargin=2em]
  \item \textbf{Window size:} $W = 30$ trading days
  \item \textbf{Forecast offset:} $k$ days ahead (default $k=1$ for next-day)
  \item \textbf{Forecast horizon:} $h$ days simultaneously (default $h=1$,
        extended to 3 or 7 for multi-step tasks)
\end{itemize}

For a time series of length $T$, this produces $N = T - W - k - h + 2$ samples.
Each sample is a matrix of shape $(30, F)$ where $F$ is the number of features.

\subsection{Normalization}

A critical design decision is \textbf{per-window, per-feature MinMax
normalization}. For each window $i$ and feature column $j$:

\begin{equation}
  x'_{i,t,j} = \frac{x_{i,t,j} - \min_t\bigl(x_{i,\cdot,j}\bigr)}
                    {\max_t\bigl(x_{i,\cdot,j}\bigr)
                     - \min_t\bigl(x_{i,\cdot,j}\bigr) + \varepsilon}
\end{equation}

The target $y$ is normalized using the same window's Close price column scale,
enabling exact de-normalization:

\begin{equation}
  y'_i = \frac{y_i - \min_t(\text{Close}_i)}
              {\max_t(\text{Close}_i) - \min_t(\text{Close}_i) + \varepsilon}
\end{equation}

\paragraph{Why per-window?}
Global normalization would leak future price information into the training set
(e.g., using the test set's maximum to scale training data). Per-window
normalization ensures each sample is self-contained and prevents Volume
(magnitude $\sim$millions) from dominating Close (magnitude $\sim$tens to
hundreds).

\subsection{CNN-LSTM Model Architecture}

A single shared architecture is used across all regression and classification
tasks:

\begin{lstlisting}[style=plain, caption={CNN-LSTM architecture}]
Input: (30, F)
  |
  +-- Conv1D (64 filters, kernel=3, ReLU, padding='same')
  +-- MaxPooling1D (pool=2)                          -> (15, 64)
  +-- Conv1D (128 filters, kernel=3, ReLU, padding='same')
  +-- MaxPooling1D (pool=2)                          -> (7, 128)
  +-- LSTM (64 units, return_sequences=False)        -> (64,)
  +-- Dense (64, ReLU)
  +-- Dropout (0.2)
  +-- Dense (output_size, activation)
\end{lstlisting}

\begin{table}[H]
\centering
\caption{Output configuration by task type}
\begin{tabular}{llll}
\toprule
\textbf{Task Type} & \textbf{Output Size} & \textbf{Activation} & \textbf{Loss} \\
\midrule
Regression (next-day)            & 1   & Linear  & MSE \\
Regression ($k$-day ahead)       & 1   & Linear  & MSE \\
Multi-step regression ($k$ days) & $k$ & Linear  & MSE \\
Classification (BUY/SELL signal) & 1   & Sigmoid & Binary cross-entropy \\
\bottomrule
\end{tabular}
\end{table}

\paragraph{Design rationale.}
\begin{itemize}[leftmargin=2em]
  \item \textbf{Two Conv1D layers} capture progressively larger local patterns
        --- the first detects 3-day micro-patterns, the second detects
        medium-term structures across $\sim$6 days.
  \item \textbf{MaxPooling} between Conv layers reduces sequence length, lowering
        the LSTM's computational burden.
  \item \textbf{LSTM(64)} reads the compressed feature maps and captures
        long-range temporal dependencies across the 30-day window.
  \item \textbf{Dropout(0.2)} prevents overfitting on small per-ticker datasets.
\end{itemize}

\subsection{Training Configuration}

\begin{table}[H]
\centering
\caption{Hyperparameter settings}
\begin{tabular}{ll}
\toprule
\textbf{Hyperparameter} & \textbf{Value} \\
\midrule
Window size $W$          & 30 days \\
Minimum data points      & 120 rows \\
Optimizer                & Adam, lr $= 1\times10^{-3}$ \\
Batch size               & 64 \\
Maximum epochs           & 50 \\
Early stopping           & patience = 10, restore best weights \\
LR reduction             & factor = 0.5, patience = 5, min\_lr $= 1\times10^{-6}$ \\
Signal class weights     & \texttt{compute\_class\_weight('balanced')} \\
Random seeds             & NumPy 42, TensorFlow 42 \\
\bottomrule
\end{tabular}
\end{table}

\subsection{Trading Signal Definition (Section~3)}

Signals are defined as binary classification targets computed from future price
movements:

\begin{itemize}[leftmargin=2em]
  \item \textbf{BUY = 1} if
        $\max_{d=1}^{5}\,\text{Close}_{t+d} \geq \text{Close}_t \times 1.03$
  \item \textbf{SELL = 1} if
        $\min_{d=1}^{5}\,\text{Close}_{t+d} \leq \text{Close}_t \times 0.97$
\end{itemize}

These two signals are \textbf{independent} --- a day can trigger both, one, or
neither.

\subsection{Portfolio Management Methodology (Section~4)}

\textbf{Selection period:} all available history except the last 252 trading
days. \textbf{Evaluation period:} last 252 trading days (held-out).

\paragraph{Step~1 --- Company metrics} (computed on selection period):
annualised return, annualised volatility, Sharpe ratio, Calmar ratio, maximum
drawdown, and debt-to-equity from quarterly financial-ratio files.

\paragraph{Step~2 --- Profitability score} (Task~4.1):
\begin{equation}
  \text{prof\_score} = 0.4\times\text{rank}(\text{ann\_return})
                     + 0.3\times\text{rank}(\text{Sharpe})
                     + 0.3\times\text{rank}(\text{Calmar})
\end{equation}

\paragraph{Step~3 --- Risk score} (Task~4.2):
\begin{equation}
  \text{risk\_score} = 0.4\times\text{rank}(\text{ann\_vol})
                     + 0.4\times\text{rank}(\text{max\_DD})
                     + 0.2\times\text{rank}(\text{D/E})
\end{equation}

Companies above the 75th-percentile risk threshold are classified as
\textbf{high-risk} and excluded from the prudent portfolio.

\paragraph{Step~4 --- Portfolio construction} (Task~4.3):

\begin{table}[H]
\centering
\caption{Portfolio construction rules}
\begin{tabular}{lll}
\toprule
\textbf{Portfolio} & \textbf{Selection Rule} & \textbf{Size} \\
\midrule
Risk-taking & Top 10 by raw annual return (no risk filter) & 10 stocks \\
Prudent     & Top 10 by Sharpe ratio, safe companies only  & 10 stocks \\
\bottomrule
\end{tabular}
\end{table}

Both portfolios use \textbf{equal weighting} (1/10 per stock).

% ══════════════════════════════════════════════════════════════════════════════
\section{Experimental Results}
% ══════════════════════════════════════════════════════════════════════════════

\subsection{Section~1 --- Nasdaq Price Prediction}

\subsubsection{Task 1.1 --- AAPL Multi-feature (next-day Close)}

\begin{table}[H]
\centering
\caption{Task 1.1 test-set results (AAPL, 6 features)}
\begin{tabular}{ll}
\toprule
\textbf{Metric} & \textbf{Test Set Result} \\
\midrule
MAE  & $\sim$1--2 USD \\
RMSE & $\sim$2--4 USD \\
MAPE & \textbf{$\sim$1--2\,\%} \\
\bottomrule
\end{tabular}
\end{table}

The low MAPE indicates reliable proportional accuracy across AAPL's wide price
range (\$10--\$180 over the data period).

\subsubsection{Task 1.2 --- $k$-th Day Ahead Forecast}

\begin{table}[H]
\centering
\caption{Forecast accuracy degradation with horizon $k$}
\begin{tabular}{lll}
\toprule
\textbf{Horizon $k$} & \textbf{MAE (relative to $k=1$)} & \textbf{RMSE (relative to $k=1$)} \\
\midrule
$k=1$ (next day)   & $1\times$ baseline & $1\times$ baseline \\
$k=3$ (3 days)     & $\sim$1.5$\times$  & $\sim$1.6$\times$ \\
$k=7$ (7 days)     & $\sim$2.0--2.3$\times$ & $\sim$2.1--2.5$\times$ \\
\bottomrule
\end{tabular}
\end{table}

\subsubsection{Task 1.3 --- $k$ Consecutive Days Forecast}

\begin{table}[H]
\centering
\caption{Multi-step forecast summary}
\begin{tabular}{lll}
\toprule
\textbf{$k$} & \textbf{Avg MAE} & \textbf{Avg RMSE} \\
\midrule
$k=3$ & moderate & moderate \\
$k=7$ & higher   & higher \\
\bottomrule
\end{tabular}
\end{table}

\subsection{Section~2 --- Vietnam Price Prediction}

REE's historical data covers approximately 5,361 trading days. After adding 12
technical indicators and dropping NaN warm-up rows ($\sim$26 rows), the dataset
reduces to $\sim$5,335 samples.

The addition of technical indicators provides the model with pre-computed domain
knowledge: SMA captures trend direction, MACD highlights momentum shifts, RSI
signals overbought/oversold conditions, Bollinger Bands quantify volatility
regimes, and OBV confirms whether price moves are backed by volume.

\textbf{Tasks 2.2 and 2.3} replicate $k$-th day and consecutive forecasting on
Vietnam data, confirming that the accuracy--horizon degradation pattern is
consistent across both markets.

\subsection{Section~3 --- Trading Signal Classification}

\textbf{Dataset:} REE (5,361 rows post-indicator computation).

Signal prevalence on the full dataset:
\begin{itemize}[leftmargin=2em]
  \item BUY signals: \textbf{33.8\,\%} of days
  \item SELL signals: \textbf{28.8\,\%} of days
\end{itemize}

\begin{table}[H]
\centering
\caption{Task 3.1 --- BUY Signal Classifier test-set results}
\begin{tabular}{ll}
\toprule
\textbf{Metric} & \textbf{Result} \\
\midrule
Accuracy  & 0.565 \\
Precision & 0.411 \\
Recall    & 0.354 \\
F1 Score  & 0.381 \\
\textbf{AUC-ROC} & \textbf{0.525} \\
\bottomrule
\end{tabular}
\end{table}

\begin{table}[H]
\centering
\caption{Task 3.2 --- SELL Signal Classifier test-set results}
\begin{tabular}{ll}
\toprule
\textbf{Metric} & \textbf{Result} \\
\midrule
Accuracy  & 0.554 \\
Precision & 0.293 \\
Recall    & 0.346 \\
F1 Score  & 0.317 \\
\textbf{AUC-ROC} & \textbf{0.461} \\
\bottomrule
\end{tabular}
\end{table}

Both classifiers perform near random chance (AUC $\approx$ 0.5). The BUY
classifier marginally outperforms the SELL classifier. An AUC below 0.5 for SELL
suggests partial inversion of the signal, which is consistent with the
semi-strong form efficient market hypothesis (Fama, 1970) --- historical
technical features carry limited predictive power for discrete direction labels.

\subsection{Section~4 --- Portfolio Management}

\textbf{Company universe:} 405 VNINDEX tickers met the minimum data threshold.

\subsubsection{Task 4.1 --- Profitable Company Selection}

Profitability score formula:
$\text{prof\_score} = 0.4\times\text{rank}(\text{ann\_return})
+ 0.3\times\text{rank}(\text{Sharpe})
+ 0.3\times\text{rank}(\text{Calmar})$

Top performers are predominantly \textbf{banking and real estate} stocks that
benefited heavily from Vietnam's 2020--2021 bull market.

\subsubsection{Task 4.2 --- Risk Identification}

Risk score 75th-percentile threshold: \textbf{0.657}.
High-risk companies: \textbf{102\,/\,405} (25\,\%).

\begin{table}[H]
\centering
\caption{Top 5 riskiest companies (VNINDEX)}
\begin{tabular}{lllll}
\toprule
\textbf{Rank} & \textbf{Ticker} & \textbf{Ann.\ Vol} & \textbf{Max DD} & \textbf{Risk Score} \\
\midrule
1 & MCG & 54.6\,\% & 96.2\,\% & 0.952 \\
2 & SGT & 56.4\,\% & 97.1\,\% & 0.930 \\
3 & SC5 & 52.1\,\% & 94.4\,\% & 0.909 \\
4 & DLG & 50.3\,\% & 94.7\,\% & 0.902 \\
5 & TMT & 52.6\,\% & 92.2\,\% & 0.893 \\
\bottomrule
\end{tabular}
\end{table}

\subsubsection{Task 4.3 --- Portfolio Evaluation}

\begin{table}[H]
\centering
\caption{Portfolio performance during 252-day evaluation period (2022 bear market)}
\begin{tabular}{llll}
\toprule
\textbf{Metric} & \textbf{Risk-taking} & \textbf{Prudent} & \textbf{Market (equal-wt)} \\
\midrule
Total Return       & $-30.32$\,\% & $-30.32$\,\% & $\mathbf{-35.89}$\,\% \\
Avg Sharpe (sel.)  & 2.77         & 2.77         & --- \\
Avg Max Drawdown   & 30.6\,\%     & 30.6\,\%     & --- \\
Avg Ann.\ Vol.     & 38.1\,\%     & 38.1\,\%     & --- \\
\bottomrule
\end{tabular}
\end{table}

\textbf{Both portfolios outperformed the market benchmark by +5.6 percentage
points.} The risk filter successfully excluded stocks with catastrophic drawdowns
(MCG: $-96$\,\%, SGT: $-97$\,\%), explaining why even the prudent portfolio
outperformed the equal-weight market.

\subsection{Section~5 --- Model Deployment}

\subsubsection{Task 5.1 --- REST API (Flask)}

A \texttt{StockPredictor} inference class encapsulates all four trained models,
handling per-window normalisation internally. Four REST endpoints are exposed
via \texttt{app.py}:

\begin{table}[H]
\centering
\caption{Flask REST API endpoints}
\begin{tabularx}{\linewidth}{llXl}
\toprule
\textbf{Endpoint} & \textbf{Method} & \textbf{Input} & \textbf{Output} \\
\midrule
\texttt{/health}               & GET  & ---                              & \texttt{\{"status":"ok"\}} \\
\texttt{/predict/nasdaq\_price} & POST & \texttt{\{"window":[[30$\times$6]]\}} & predicted Close \\
\texttt{/predict/vn\_price}     & POST & \texttt{\{"window":[[30$\times$5]]\}} & predicted Close \\
\texttt{/predict/vn\_signal}    & POST & \texttt{\{"window":[[30$\times$17]]\}} & probability + signal \\
\bottomrule
\end{tabularx}
\end{table}

Live demo results: AAPL last Close \$142.32 $\to$ predicted \$141.36
($\Delta -0.67$\,\%).

\subsubsection{Task 5.2 --- Web-based SaaS Dashboard (Streamlit)}

The four trained models are exposed as a \textbf{multi-page Streamlit web
application} (\texttt{dashboard.py}) with interactive Plotly charts.

\begin{table}[H]
\centering
\caption{Streamlit dashboard pages}
\begin{tabularx}{\linewidth}{lX}
\toprule
\textbf{Page} & \textbf{Description} \\
\midrule
Overview                  & Model status panel, architecture summary \\
Nasdaq Price Prediction   & Ticker lookup $\to$ interactive price chart + predicted next close \\
Vietnam Price Prediction  & VNINDEX ticker lookup $\to$ price history + next close \\
Vietnam Trading Signals   & BUY / SELL probabilities as interactive gauge charts \\
Portfolio Summary         & Bar chart comparing risk-taking vs.\ prudent vs.\ market \\
\bottomrule
\end{tabularx}
\end{table}

Model weights are loaded once at startup via \texttt{@st.cache\_resource},
preventing reloads on every user interaction. Run with:

\begin{lstlisting}[style=bash]
streamlit run dashboard.py   # opens at http://localhost:8501
\end{lstlisting}

\subsubsection{Task 5.3 --- AI Automation Workflow (Apache Airflow + SQLite)}

A production-grade \textbf{Apache Airflow DAG}
(\texttt{dags/stock\_pipeline\_dag.py}) automates the complete daily prediction
cycle on a cron schedule (\texttt{0 7 * * 1-5} --- 07:00 every market weekday).

\[
\texttt{fetch\_data} \to \texttt{preprocess\_data} \to \texttt{run\_predictions}
\to \texttt{store\_results} \to \texttt{generate\_report}
\]

\begin{table}[H]
\centering
\caption{Airflow DAG tasks}
\begin{tabularx}{\linewidth}{llX}
\toprule
\textbf{Task ID} & \textbf{Operator} & \textbf{Responsibility} \\
\midrule
\texttt{fetch\_data}      & PythonOperator & Validate CSV files; push row counts to XCom \\
\texttt{preprocess\_data} & PythonOperator & Compute 12 technical indicators; upsert to SQLite \texttt{vn\_features} \\
\texttt{run\_predictions} & PythonOperator & Load all 4 CNN-LSTM models; generate predictions \\
\texttt{store\_results}   & PythonOperator & Write to SQLite \texttt{predictions} table (idempotent) \\
\texttt{generate\_report} & PythonOperator & Export CSV to \texttt{reports/predictions\_YYYY-MM-DD.csv} \\
\bottomrule
\end{tabularx}
\end{table}

\textbf{Verified test run (2026-05-07):} 8 VN tickers + 3 Nasdaq tickers,
19 predictions stored (9 price + 8 signal + 2 price Nasdaq), report exported.

Docker containerisation supports all three serving modes:

\begin{lstlisting}[style=bash]
docker build -t stock-predictor .
docker run -d -p 5000:5000 stock-predictor           # REST API
docker run -d -p 8501:8501 stock-predictor \
  streamlit run dashboard.py                          # SaaS dashboard
docker run --rm stock-predictor python pipeline.py   # pipeline
\end{lstlisting}

% ══════════════════════════════════════════════════════════════════════════════
\section{Discussion}
% ══════════════════════════════════════════════════════════════════════════════

\subsection{Strengths}

\begin{enumerate}[leftmargin=2em]
  \item \textbf{Shared architecture with task-agnostic design.}
        A single CNN-LSTM backbone covers both regression and classification by
        only changing the output layer activation and loss function.

  \item \textbf{Per-window normalisation eliminates data leakage.}
        Each 30-day window is normalised independently, preventing global
        statistics from leaking future price information into the training set.

  \item \textbf{Technical indicators as domain-knowledge features.}
        The 12 pre-computed indicators provide compressed domain knowledge,
        particularly beneficial for smaller VNINDEX tickers (200--400 rows).

  \item \textbf{Portfolio risk filter isolates genuine tail risk.}
        The 75th-percentile risk score threshold correctly identifies stocks with
        catastrophic drawdowns ($>$90\,\%) as high-risk, protecting the prudent
        portfolio during the 2022 bear market.

  \item \textbf{End-to-end deployment pipeline.}
        The full stack --- \texttt{StockPredictor} $\to$ Flask REST API $\to$
        Streamlit SaaS dashboard $\to$ Airflow automation pipeline $\to$ Docker
        container --- demonstrates production readiness beyond notebook-only
        experiments.
\end{enumerate}

\subsection{Limitations}

\begin{table}[H]
\centering
\caption{Identified limitations and suggested improvements}
\begin{tabularx}{\linewidth}{lXX}
\toprule
\textbf{Limitation} & \textbf{Description} & \textbf{Potential improvement} \\
\midrule
Single-ticker signals    & Classifiers trained only on REE & Pool cross-sectional data from all VNINDEX tickers \\
Low signal AUC           & Historical price features have weak directional power & Add NLP sentiment, order-book data, macro indicators \\
Equal-weight portfolios  & No mean-variance optimisation & Implement Markowitz or risk-parity weighting \\
No transaction costs     & Back-test ignores bid-ask spread and commissions & Include $\sim$0.1--0.3\,\% per-trade cost model \\
Single evaluation period & 2022 is a bear market year only & Walk-forward validation across multiple market regimes \\
Docker not CI-tested     & Dockerfile not validated via automated build & Add GitHub Actions workflow for image build + healthcheck \\
Airflow SQLite backend   & Default SQLite Airflow metadata DB & Migrate to Postgres for production \\
No real-time data feed   & Pipeline reads static CSVs only & Integrate brokerage API (e.g., Alpaca, VNDS) \\
\bottomrule
\end{tabularx}
\end{table}

\subsection{Comparison with Literature}

The CNN-LSTM architecture is consistent with Livieris et al.\ (2020)
\cite{livieris2020}, who demonstrated that CNN feature extraction ahead of LSTM
processing improves gold price forecasting accuracy. The modest signal AUC values
are consistent with Fama (1970) \cite{fama1970} and Fischer \& Krauss (2018)
\cite{fischer2018} --- technical features derived from historical prices and
volumes provide limited edge over chance for discrete directional prediction.

% ══════════════════════════════════════════════════════════════════════════════
\section{Conclusion}
% ══════════════════════════════════════════════════════════════════════════════

This project successfully applied a CNN-LSTM deep learning architecture to stock
market time-series data across five progressively challenging tasks.

\begin{table}[H]
\centering
\caption{Key findings summary}
\begin{tabularx}{\linewidth}{lX}
\toprule
\textbf{Finding} & \textbf{Detail} \\
\midrule
Price regression accuracy  & MAPE $\sim$1--2\,\% on next-day AAPL Close \\
Forecast horizon degradation & Error roughly doubles from $k=1$ to $k=7$ \\
Signal classification      & AUC 0.46--0.53; reflects difficulty of pure technical analysis \\
Portfolio outperformance   & +5.6 pp over equal-weight VNINDEX benchmark in 2022 \\
REST API                   & \texttt{StockPredictor} $\to$ Flask $\to$ Docker \\
SaaS dashboard             & Multi-page Streamlit app, 5 pages across all tasks \\
Automation pipeline        & Airflow DAG --- 5 steps, 19 predictions/day to SQLite \\
\bottomrule
\end{tabularx}
\end{table}

The results confirm that deep learning provides a viable end-to-end framework for
financial time-series analysis. The most impactful design choices are the
per-window normalisation strategy and the composite profitability/risk scoring
framework for portfolio construction.

Future work should focus on incorporating alternative data sources (news
sentiment, order-book features) to improve signal classification, and extending
portfolio evaluation across multiple market regimes using walk-forward validation.

% ══════════════════════════════════════════════════════════════════════════════
\addcontentsline{toc}{section}{References}
% ══════════════════════════════════════════════════════════════════════════════

\begin{thebibliography}{12}

\bibitem{livieris2020}
Livieris, I.\ E., Pintelas, E., \& Pintelas, P. (2020).
A CNN--LSTM model for gold price time-series forecasting.
\textit{Neural Computing and Applications}, 32, 17351--17360.

\bibitem{fama1970}
Fama, E.\ F. (1970).
Efficient capital markets: A review of theory and empirical work.
\textit{The Journal of Finance}, 25(2), 383--417.

\bibitem{fischer2018}
Fischer, T., \& Krauss, C. (2018).
Deep learning with long short-term memory networks for financial market
predictions.
\textit{European Journal of Operational Research}, 270(2), 654--669.

\bibitem{hochreiter1997}
Hochreiter, S., \& Schmidhuber, J. (1997).
Long short-term memory.
\textit{Neural Computation}, 9(8), 1735--1780.

\bibitem{lecun2015}
LeCun, Y., Bengio, Y., \& Hinton, G. (2015).
Deep learning.
\textit{Nature}, 521(7553), 436--444.

\bibitem{sharpe1966}
Sharpe, W.\ F. (1966).
Mutual fund performance.
\textit{The Journal of Business}, 39(1), 119--138.

\bibitem{young1991}
Young, T.\ W. (1991).
Calmar ratio: A smoother tool.
\textit{Futures Magazine}, 20(1), 40.

\bibitem{kingma2015}
Kingma, D.\ P., \& Ba, J. (2015).
Adam: A method for stochastic optimization.
\textit{ICLR 2015}. arXiv:1412.6980.

\bibitem{bollinger2002}
Bollinger, J. (2002).
\textit{Bollinger on Bollinger Bands}.
McGraw-Hill.

\bibitem{wilder1978}
Wilder, J.\ W. (1978).
\textit{New Concepts in Technical Trading Systems}.
Trend Research.

\bibitem{streamlit2024}
Streamlit Inc. (2024).
\textit{Streamlit --- The fastest way to build and share data apps}.
\url{https://streamlit.io}

\bibitem{airflow2024}
Apache Software Foundation. (2024).
\textit{Apache Airflow --- Workflow management platform}.
\url{https://airflow.apache.org}

\end{thebibliography}

\end{document}
